%定义知识点拓展盒子knowledgebox----------------------------------
\newtcolorbox{knowledgebox}[2][]
  {enhanced,
    before skip=2mm,
    after skip=2mm,
    breakable,
    colback=c4,
    colframe=c1,%边框颜色
    boxrule=0.5mm,
    attach boxed title to top left={xshift=1cm,yshift*=1mm-\tcboxedtitleheight},
    fonttitle=\bfseries,
    % varwidth boxed title*=-3cm,
	boxed title style={frame code={
                    \path[fill=c1]
                    ([yshift=-1mm,xshift=-1mm]frame.north west)
                    arc[start angle=0,end angle=180,radius=1mm]
                    ([yshift=-1mm,xshift=1mm]frame.north east)
                    arc[start angle=180,end angle=0,radius=1mm];
                    \path[left color=c1,right color=c1,
                        middle color=c1]
                    ([xshift=-2mm]frame.north west) -- ([xshift=2mm]frame.north east)
                    [rounded corners=1mm]-- ([xshift=1mm,yshift=-1mm]frame.north east)
                    -- (frame.south east) -- (frame.south west)
                    -- ([xshift=-1mm,yshift=-1mm]frame.north west)
                    [sharp corners]-- cycle;
                },interior engine=empty,
        },
    title=\large \( \textbf{解释说明： #2} \),#1}

%定义解题点拨盒子skillbox----------------------------------
\newtcolorbox{skillbox}[2][]
  {enhanced,
  breakable,
  sharp corners,
  attach boxed title to top left={
    yshifttext=-1mm
  },
  colback=c4,
  colframe=c1,
  boxrule=0.5mm,
  fonttitle=\bfseries,
  coltitle=white,
  boxed title style={
    sharp corners,
    size=small,
    colback=c1,
    colframe=c1,
  },
  title=\large \( \textbf{解题点拨： #2} \),
  #1}



%定义考查方式盒子pointbox----------------------------------
\newtcolorbox{pointbox}[2][]
  {enhanced,
  breakable,
  colback = c4,
  frame hidden,
  borderline west = {2pt}{0pt}{c1},
  coltitle = c0,
  fonttitle = \bfseries\sffamily,
  boxrule = 0sp,
  sharp corners,
  detach title,
  before upper = \tcbtitle\par\smallskip,
  description font = \mdseries,
  separator sign none,
  segmentation style={solid, c1},
	title=\large \( \textbf{考查方式： #2} \),#1}

% 备用盒子1boxa----------------------------------
\makeatletter
\newtcolorbox{boxa}[1]{%
    enhanced,
    breakable,
    colback=c4,
    colframe=c0,
    attach boxed title to top left={yshift*=-\tcboxedtitleheight},
    fonttitle=\bfseries,
    title={\large \( \textbf{备用盒子： #1} \)},
    boxed title size=title,
    boxed title style={%
            sharp corners,
            rounded corners=northwest,
            colback=c0,
            boxrule=0pt,
        },
    underlay boxed title={%
            \path[fill=c0] (title.south west)--(title.south east)
            to[out=0, in=180] ([xshift=5mm]title.east)--
            (title.center-|frame.east)
            [rounded corners=\kvtcb@arc] |-
            (frame.north) -| cycle;
        }
}
\makeatletter

%例题盒子定义example----------------------------------
\newtcolorbox{example}[1]{
    colback = c4,
    breakable,
    colframe = c1,
    coltitle = c0,
    title={\large \( \textbf{Example： #1} \)},
    boxrule = 1pt,
    sharp corners,
    detach title,
    before upper=\tcbtitle\par\smallskip,
    fonttitle = \bfseries,
    description font = \mdseries
}

%定义注释盒子note----------------------------------
\newtcolorbox{note}[1][]{%
    enhanced jigsaw,
    colback=c4,
    colframe=c1,
    size=small,
    boxrule=1pt,
    title=\textbf{Note},
    halign title=flush center,
    coltitle=black,
    breakable,
    drop shadow=c2,
    attach boxed title to top left={xshift=1cm,yshift=-\tcboxedtitleheight/2,yshifttext=-\tcboxedtitleheight/2},
    minipage boxed title=1.5cm,
    boxed title style={%
        colback=white,
        size=fbox,
        boxrule=1pt,
        boxsep=2pt,
        underlay={%
            \coordinate (dotA) at ($(interior.west) + (-0.5pt,0)$);
            \coordinate (dotB) at ($(interior.east) + (0.5pt,0)$);
            \begin{scope}
                \clip (interior.north west) rectangle ([xshift=3ex]interior.east);
                \filldraw [white, blur shadow={shadow opacity=60, shadow yshift=-.75ex}, rounded corners=2pt] (interior.north west) rectangle (interior.south east);
            \end{scope}
            \begin{scope}[c1]
                \fill (dotA) circle (2pt);
                \fill (dotB) circle (2pt);
            \end{scope}
        },
        },
        #1,
}

%定义考点习题标题及分割线\et----------------------------------
\newcommand{\et}[0]{
  \color{c1}
  \subsubsection*{\large 考点习题}
  \hrule
  \vspace{1em}
  \color{black}
}

%初始盒子1----------------------------------------
\newtcolorbox{warningbox}{
	colback=green!10,
	colframe=olive!50!black,
	leftrule=0.5mm,
	rightrule=0mm,
	toprule=0mm,
	bottomrule=0mm,
	arc=2pt,
	outer arc=2pt
}

% 初始盒子2：定义定理框样式------------------------------
\newtcolorbox{theorembox}{
	colback=gray!10,
	colframe=olive!50!black,
    leftrule=0.5mm,
    rightrule=0mm,
    toprule=0mm,
    bottomrule=0mm,
	arc=2pt,
	outer arc=2pt
}